\PassOptionsToPackage{sols}{handout}
\documentclass[main]{subfiles}
\setcounterrefbeforedot{chapter}{m-ws7-2}
\setcounterrefafterdot{section}{m-ws7-2}
\addtocounter{section}{-1}
\usepackage{solutions}

\begin{document}

\ifSubfilesClassLoaded{%
  \chaptermark{\chaptersix}
}{}

\section{Modeling with Differential Equations} \label{ws7-2}

\begin{problemonly}
    \begin{framed}

\textbf{Key Ideas}:

\begin{itemize}[topsep=0pt,partopsep=0pt,parsep=8pt,itemsep=0pt]

\item You should be able to write a differential equation that reasonably
  models a given scenario.  You should also be able to write down an
  initial condition, if the scenario includes one.

\item You should understand how continuously compounded interest can be
  expressed using a differential equation.

\end{itemize}
\end{framed}
\end{problemonly}

Let's take a moment to examine the difference between the differential equations 
$\dfrac{dy}{dt}=y$ and $\dfrac{dy}{dt}=t$.

\begin{itemize}[topsep=0pt]
    \item $\dfrac{dy}{dt}=y$ is a differential equation of the form
          $\dfrac{dy}{dt}=f(y)$. The rate of change of $y$ is a function
          of \mbox{$y$ itself.}
    \item $\dfrac{dy}{dt}=t$ is a differential equation of the form
          $\dfrac{dy}{dt}=f(t)$. The rate of change of $y$ is a function
          of $t$.
\end{itemize}

The general solution of $\dfrac{dy}{dt}=y$ is $y(t)=\cfitb{1in}{Ce^t}$.

The general solution of $\dfrac{dy}{dt}=t$ is $y(t)=\cfitb{1in}{\frac{1}{2}t^2+C}$.
    \begin{enumerate}

\item
\note{Students struggle with the difference between a differential equation
and its solution when they first see differential equations. These next
two questions are to help them tackle this head on.}
    \begin{enumerate}
        \item Last class we talked about a population of bacteria, denoted
        $P(t)$, that was growing at a rate proportional to itself. We found the
        proportionality constant was $0.2$. We found a differential equation
        modeling the situation, and a solution corresponding to an initial
        condition.
        \begin{center}
            \begin{tabular}{ll}
            Differential equation: \cfitb{1in}{$\dfrac{dP}{dt}=0.2P$} &
            General solution: \cfitb{1in}{$P(t)=Ce^{0.2t}$} \\
            & \\
            Initial condition: \underline{\parbox[t]{1in}{\centering$P(0)=200$}} &
            Corresponding particular solution: \cfitb{1in}{$P(t)=200e^{0.2t}$}
            \end{tabular}
        \end{center}
        
  %  \pspace{\newpage}
    
    \item Ebony puts $\$2000$ in a bank account earning a nominal annual interest
    rate of $3\%$ compounded continuously. She makes no further deposits or
    withdrawals on this particular account. Let $M(t)$ be the amount
    of money in her account at time $t$. (Note that money in an interest-earning
    account grows at a rate proportional to itself.)

    Fill in the blanks from the options below. \vspace{-12pt} 
    \begin{center}
            \begin{tabular}{ll}
            Differential equation: \cfitb{1in}{$\dfrac{dM}{dt}=0.03M$} &
            General solution: \cfitb{1in}{$M(t)=Ce^{0.03t}$} \\
            & \\
            Initial condition: 
            \cfitb{1in}{$M(0)=2000$}&
            Corresponding particular solution: \cfitb{1.1in}{$M(t)=2000e^{0.03t}$}
            \end{tabular}
        \end{center}
    \begin{multicols}{2}
        \begin{enumerate}[topsep=0pt]
        \item $M(0)=2000$
        \item $M(t)=2000e^{0.03t}$
        \item $M(t)=2000(1+0.03)^t$
        \item $\frac{dM}{dt}=2000e^{0.03t}$
        \item $M(t)=Ce^{0.03t}$
        \item $\frac{dM}{dt}=Ce^{0.03t}$
        \columnbreak
        \item $\frac{dM}{dt}=0.03$
        \item $\frac{dM}{dt}=1.03$
        \item $\frac{dM}{dt}=0.03t$
        \item $\frac{dM}{dt}=0.03M$
        \item $\frac{dM}{dt}=1.03M$
        \end{enumerate}
    \end{multicols}
    \end{enumerate}

\end{enumerate}
\vfill
  \begin{framed}
    \textbf{Key Idea 1.} Saying that two quantities $A$ and $B$ are
    \uline{proportional} to each other means that one is a constant
    multiple of the other.  That is,
    \[ A = k B \]
    for some constant $k$.
  \end{framed}


\begin{enumerate}[resume]
\item 
  \begin{problem}[\vfill]
    The population of a certain country increases at a rate proportional
    to the population size.  If $P(t)$ is the population size at time $t$,
    write a differential equation that models the situation. 
  \end{problem}

  \begin{solution}
    We are told that
    \[ \underbrace{\text{The rate at which the population
        increases}}_{\frac{dP}{dt}} \text{ } \underbrace{\text{is
        proportional to}}_{= k \text{ times}} \text{ }
    \underbrace{\text{the population size}}_{P}.\] So, our differential
    equation is $\boxed{\frac{dP}{dt} = k P}$.  (Since the population is
    increasing, $\frac{dP}{dt}$ is positive.  The population $P$ is also
    positive, so the constant $k$ must be positive.)
  \end{solution}

\item 
  One way that psychologists study how people learn is by asking them to
  learn a list of nonsense words.\footnote{This is meant to be a ``pure''
    experiment, in that the learners won't have any prior knowledge that
    would help them with the task.}  Suppose that $L(t)$ is the percent of
  such a list that a particular person has learned at time $t$, where $t$
  is measured in days.
  \begin{enumerate}
  \item

    \note{You might like to have students talk about whether this model
      seems reasonable.}

    \begin{problem}[\vfill]
      A simple model for learning is to assume that learning happens at a
      rate proportional to the percent of the list left to be learned.
      Write a differential equation that expresses this.  Can you write an
      initial condition?
    \end{problem}

    \begin{solution}
      In this model, we assume that the rate at which learning happens
      ($\frac{dL}{dt}$) is proportional to the percent of the list left to
      be learned ($100 - L$).  That is, $\boxed{\frac{dL}{dt} = k (100 -
        L)}$. 
    \end{solution}

  \item
    \begin{problem}[\stretch{1.5}]
      A student who has learned a fifth of the list already is currently
      learning at a rate of 10\% of the list per day.  Incorporate this
      into your model.
    \end{problem}

    \begin{solution}
      At this particular moment, $L = 20$ (the student has learned 20\% of
      the list) and $\frac{dL}{dt} = 10$ (the student is learning at a
      rate of 10\% per day).  If we plug this into the differential
      equation $\frac{dL}{dt} = k(100 - L)$, then we find that $10 = k
      (100 - 20)$, so we are able to solve for $k$: it's $\frac{10}{80} =
      \frac{1}{8}$.  Thus, $\boxed{\frac{dL}{dt} = \frac{1}{8} (100 -
        L)}$.
    \end{solution}

  \end{enumerate}

\item 

  \note{This can also be a basic model for the spread of a disease, which
    students usually find really interesting.}

  A juicy rumor is spreading through Harvard's freshman class, which
  has 1600 students.  The rumor spreads at a rate proportional to the
  product of the number of freshmen who have already heard it and the
  number of freshmen who have not yet heard it.  At time $t = 0$, 5
  freshmen have heard the rumor.

  \begin{enumerate}
  \item 
    \begin{problem}[\vfill]
      If $R(t)$ is the number of people who have heard the rumor at time
      $t$, write a differential equation that models the situation.  Can
      you write an initial condition?
    \end{problem}

    \begin{solution}
      We are told that the rate at which the rumor spreads
      ($\frac{dR}{dt}$) is proportional to the product of the number of
      freshmen who have already heard it ($R$) and the number of freshmen
      who have not yet heard it ($1600 - R$).  So, $\boxed{\frac{dR}{dt} =
        kR(1600-R)}$.  We are also told that $\boxed{R(0) = 5}$; this is
      the initial condition.
    \end{solution}

  \item
    \begin{problem}[\vfill]
      Lucy thinks that the rumor should instead spread at a rate
      proportional to the number of students who haven't heard the rumor
      yet.  Which model (Lucy's or the original one) do you think is more
      reasonable?  Why?
    \end{problem}

    \begin{solution}
      Lucy's model is $\frac{dR}{dt} =k(1600-R)$.

      To decide which is more reasonable, we can think about some special
      cases; if nobody has heard the rumor ($R = 0$), we'd expect
      $\frac{dR}{dt}$ to be 0 (nobody knows the rumor, so nobody can
      spread it).  This is true in the original model but not in Lucy's
      model, so the original model seems to make more sense. 
    \end{solution}

    \item
    \begin{problem}[\vfill]
        Could a similar model be used for the spread of an infectious disease?
        Why or why not? (Or in what circumstances could it be used?)
    \end{problem}

    \begin{solution}
        In epidemiology, a very simple model for the spread of an infectious
        disease could be $\dfrac{dI}{dt}=kI(P-I)$, where $k$ is a constant, $I$
        is the number of infected individuals, and $P$ is the population size.
        In practice, epidemiologists use more complicated models that also take
        into account the number of susceptible individuals $S$ and the number
        of recovered individuals $R$. (This is called the SIR model.)

        This model can also be used in ecology to model population growth, e.g.
        $\dfrac{dP}{dt}=\dfrac{r}{K}P(K-P)$, where $P$ is population size,
        $r$ is growth rate, and $K$ is the carrying capacity (the maximum
        population an ecological system can support.) It can also be used
        in chemistry to model reaction rates.
    \end{solution}

  \end{enumerate}

\end{enumerate}


\probonly{\newpage}

  \begin{framed}
    \textbf{Key Idea 2.}  Differential equations can be used to model more
    complicated situations.  When we want to do this, a key equation is
    \begin{align*}
      \text{overall rate of change} &= \text{(rate in)} - \text{(rate out)} \\
      \textit{or} \hspace{0.5in} \text{overall rate of change} &= \text{(rate of
        increase)} - \text{(rate of decrease)}
    \end{align*}
  \end{framed}

   % \pspace{\newpage}

\begin{enumerate}[resume]
\item 

  \note{Students often have trouble putting all of the information
    together.  The first principle they need to recognize is that the rate
    of change is (rate in) $-$ (rate out).  Some students will want the
    rate in to involve $t$ somehow.  For instance, one common mistake it
    to say the rate in is $10t$; what they seem to be thinking is that,
    after $t$ minutes, the \emph{amount} that has been administered is
    $10t$.  It may help to emphasize that the rate is a constant rate so
    should not depend on $t$. Thinking about the units of $\frac{dM}{dt}$
    may also help.}



  \begin{problem}[1in]
    A drug is being administered to a patient at a constant rate of $10$
    mg/hr.  The patient metabolizes and eliminates the drug at a rate
    proportional to the amount in his body. Let $M = M(t)$ be the amount
    (in mg) of medicine in the patient's body at time $t$, where $t$ is
    measured in hours.  Write a differential equation that models the
    situation.
  \end{problem}

  \begin{solution}
    The rate of change of medicine in the patient's body is equal to (rate
    in) $-$ (rate out).  The rate in of medicine is $10$ mg/hr.  The rate
    out is proportional to the amount of medicine in his body, so it's
    $kM$ for some positive constant $k$.  Therefore, the appropriate model
    is $\boxed{\frac{dM}{dt} = 10 - kM}$.
  \end{solution}

\item Estella opens a bank account with \$20,000. The account has a nominal
annual interest rate of 4\% per year compounded continuously. She withdraws
money at a rate of \$2400/year (about \$200/month.) 
\begin{enumerate}
    \item 
    \begin{problem}[1in] Write a differential equation for $M(t)$,
    the amount of money in her account at time $t$. What is the initial condition?
    \end{problem}

    \begin{solution}
        The rate of change of $M(t)$ over time, $\frac{dM}{dt}$, is equal
        to the rate at which Estella earns interest minus the rate at which
        she withdraws money. The rate at which she earns interest is proportional
        to $M(t)$, $0.04M$. She withdraws money at a constant rate of 
        2400 (in dollars per year, and we assume that she withdraws money
        continuously.) So \fbox{$\frac{dM}{dt}=0.04M-2400$.} The initial
        condition is \fbox{$M(0)=20,000$.}
    \end{solution}

    \item Which of the following gives Estella's balance at time $t$?
    \begin{enumerate}
        \item $M(t)=20,000e^{0.04t}-2400t$
        \probonly{\item $M(t)=60,000-40,000e^{0.04t}$}
        \solonly{\item \fbox{$M(t)=60,000-40,000e^{0.04t}$}}
    \end{enumerate}
    \pspace{\vfill}

    \begin{solution}
        Let's check.
        \begin{enumerate}
            \item If $M(t)=20,000e^{0.04t}-2400t$, then 
            $\frac{dM}{dt}=800e^{0.04t}-2400$. If we substitute $M(t)$ into
            the right side of the differential equation, we get 
            \begin{align*}
            0.04M-2400 &= 0.04(20,000e^{0.04t}-2400t)-2400\\
            &=800e^{0.04t}-96t-2400.
            \end{align*}
            This is not the same as $\frac{dM}{dt}$, so this is not a solution. 
            \item If $M(t)=60,000-40,000e^{0.04t}$, then
            $\frac{dM}{dt}=-1600e^{0.04t}$. If we substitute $M(t)$ into the
            right side of the differential equation, we get
            \begin{align*}
            0.04M-2400 &= 0.04(60,000-40,000e^{0.04t})-2400\\
            &=-1600e^{0.04t}.
            \end{align*}
            This is the same as $\frac{dM}{dt}$, so this is a solution.
        \end{enumerate}
    \end{solution}
\end{enumerate}


\probonly{\newpage}

\item 

  Dante borrows \$150,000 from the bank.  He pays money back at a rate of
  \$1000 per month, while the bank charges him interest at a nominal
  annual interest rate of 6\% per year compounded continuously.

  \begin{enumerate}
  \item

    \note{I've made this a monthly payment rather than yearly for two
      reasons: that's typically how mortgages are structured, and a
      monthly payment feels more continuous than a yearly one.
      Nonetheless, students may object to modeling the \$12000/year
      continuously, so you may want to point out that we very often model
      discrete phenomena continuously (e.g., when we model the number of
      people in a town at time $t$).}

    \begin{problem}[1in]
      Write a differential equation that models $B(t)$, the balance Dante
      owes the bank at time $t$, where $t$ is measured in years and $t =
      0$ is the time at which Dante borrows the money.  \emph{Be careful
        with units!} \emph{Note that for the purposes of our model, we will model
        the payments, which total to \$12,000/year, as if they were being made
        continuously.}
    \end{problem}

    \begin{solution}
      $\frac{dB}{dt}$ is the rate of change of Dante's balance; since $B$
      is in dollars and $t$ is in years, this rate should be in dollars
      per year.  This rate is equal to (rate of increase due to interest)
      $-$ (rate of decrease due to payments).  The rate of increase due to
      interest is $0.06 B$, while the rate of decrease due to payments is
      a constant \$1000 per month, or \$12,000 per year.  So,
      $\boxed{\frac{dB}{dt} = 0.06B - 12,000}$.  We also have an initial
      condition, $\boxed{B(0) = 150,000}$, since we know that Dante owes
      \$150,000 at $t = 0$.
    \end{solution}

  \item

    \note{Here, I'd like to make the point that figuring this out directly
      from the scenario is very hard; however, once we've written down the
      differential equation, it's easy to check whether a given function
      satisfies the differential equation.  What I want them to get out of
      this problem in the end is that writing down $B(t)$ directly is
      actually very hard, but it's easy to model the situation using
      differential equation.  I'll probably say a little bit about the
      fact that there are methods that allow us to find the solution of
      the differential equation methodically, although we won't cover
      those in this course.

      I'll also try to use the scenario to convince students that there is
      exactly one function that satisfies the differential equation and
      initial condition.
    }

    \begin{problem}[\vfill]
      Which of the following gives Dante's balance at time $t$? 
      \begin{enumerate}[label=\protect\maybecircle{\roman*.}]
      \plainitem $B(t) = 150,000 e^{0.06 t} - 12,000 t$
      \circleitem $B(t) = 200,000 - 50,000 e^{0.06 t}$
      \plainitem $B(t) = (150,000 - 12,000 t) e^{0.06 t}$
      \end{enumerate}
    \end{problem}

    \begin{solution}
      We simply need to determine which of the choices satisfies both the
      differential equation $\frac{dB}{dt} = 0.06 B - 12,000$ and the
      initial condition $B(0) = 150,000$.  Let's look at each one:
      \begin{enumerate}[label=\roman*.]
      \item If $B(t) = 150,000 e^{0.06 t} - 12,000 t$, then $\frac{dB}{dt}
        = 0.06 (150,000 e^{0.06 t}) - 12,000$ while $0.06 B - 12,000 =
        0.06 (150,000 e^{0.06 t} - 12,000 t) - 12,000$.  So,
        $\frac{dB}{dt} \neq 0.06 B - 12,000$.
      \item If $B(t) = 200,000 - 50,000 e^{0.06 t}$, then $\frac{dB}{dt} =
        - 0.06 (50,000 e^{0.06 t})$ while $0.06 B - 12,000 = 0.06 \left(
          200,000 - 50,000 e^{0.06 t} \right) - 12,000$, which simplifies
        to $- 0.06 (50,000 e^{0.06 t})$.  These are equal, so we at least
        have a solution of the differential equation.  We should also see
        if this function satisfies the initial condition: $B(0) = 200,000
        - 50,000 e^{0} = 150,000$.  So, this must be Dante's balance at
        time $t$.
      \item We don't need to check this one, because we already figured
        out that $B(t) = 200,000 - 50,000 e^{0.06 t}$.  (There is exactly
        one function that satisfies both the differential equation and the
        initial condition.)
      \end{enumerate}
    \end{solution}

  \end{enumerate}

\end{enumerate}

\end{document}